% !TEX program = lualatex
\documentclass[professionalfonts]{beamer}
\newcommand{\slidemode}{1}  % カラー投影モード
\usepackage{beamer_template}
\usepackage{enumerate}

\newcommand{\LectureNum}{5}
\newcommand{\LectureTitle}{微分方程式への応用}
\newcommand{\CourseName}{応用数学}
\renewcommand{\TermName}{前期}

\begin{document}

\begin{frame}{第5回　微分方程式への応用}
  \begin{block}{今回の内容}
  微分方程式への応用
  \end{block}
  \vfill\hfill{\scriptsize \textcolor{gray}{第2章 ラプラス変換　§2}}
\end{frame}

\begin{frame}{ラプラス変換による微分方程式の解法}
  \begin{block}{}
  \begin{itemize}
    \item 微分方程式と初期条件をラプラス変換
    $\Rightarrow$ $X(s)$ に関する代数方程式
    \item 代数方程式を $X(s)$ について解く
    $\Rightarrow$ $X(s)$ の分数式（部分分数分解）
    \item $X(s)$ を逆ラプラス変換
    $\Rightarrow$ 解 $x(t)$
  \end{itemize}
  \vspace{4pt}
  {\small ※ 代数方程式：微分を含まない，四則演算だけの方程式}
  \end{block}
\end{frame}

\begin{frame}{例題1}
  \begin{exampleblock}{問題}
    $dx/dt + x = e^t, x(0) = 1$ を解け\\[2pt]
  \end{exampleblock}
\end{frame}

\begin{frame}{例題1　解答 (1/2)}
  $\mathcal{L}[x(t)] = X(s)$ とおく．原関数の微分法則より\\[4pt]
  $\mathcal{L}\!\left[\dfrac{dx}{dt}\right] = sX(s) - x(0) = sX(s) - 1$\\[4pt]
  両辺をラプラス変換：$(sX(s) - 1) + X(s) = \dfrac{1}{s-1}$\\[4pt]
  $(s+1)X(s) = 1 + \dfrac{1}{s-1} = \dfrac{s}{s-1}$\\[6pt]
  $\therefore\quad X(s) = \dfrac{s}{(s+1)(s-1)}$
\end{frame}

\begin{frame}{例題1　解答 (2/2)}
  部分分数分解：$\dfrac{s}{(s+1)(s-1)} = \dfrac{A}{s-1}+\dfrac{B}{s+1}$\\[4pt]
  通分：$s = A(s+1) + B(s-1)$\\[4pt]
  $s=1$：$1=2A \Rightarrow A=\dfrac{1}{2}$，\quad $s=-1$：$-1=-2B \Rightarrow B=\dfrac{1}{2}$\\[4pt]
  $X(s) = \dfrac{1/2}{s-1} + \dfrac{1/2}{s+1}$\\[4pt]
  逆ラプラス変換：$x(t) = \dfrac{1}{2}e^t + \dfrac{1}{2}e^{-t}$
  \\[8pt]\textcolor{red}{\textbf{答}}：$x(t) = \cosh t$
\end{frame}

\begin{frame}{演習問題}
  \begin{exampleblock}{自分で解いてみよう}
  \begin{description}
    \item[問1] 次の微分方程式を解け
    $(1) dx/dt = 2x,  x(0) = 1$\\
    $(2) dx/dt + x = e^{-t},  x(0) = 0$\\
  \end{description}
  \end{exampleblock}
\end{frame}

\begin{frame}{問1　解答}
  (1) ラプラス変換：$sX(s) - 1 = 2X(s)$ \quad $\therefore (s-2)X(s) = 1$\\[4pt]
  $X(s) = \dfrac{1}{s-2}$ \quad $\therefore x(t) = e^{2t}$\\[6pt]
  (2) ラプラス変換：$(sX(s)-0) + X(s) = \dfrac{1}{s+1}$ \quad $\therefore (s+1)X(s) = \dfrac{1}{s+1}$\\[4pt]
  $X(s) = \dfrac{1}{(s+1)^{2}}$\\[4pt]
  逆変換：$x(t) = te^{-t}$
  \\[8pt]\textcolor{red}{\textbf{答}}：
  $(1)\ x(t) = e^{2t}$\\
  $(2)\ x(t) = te^{-t}$\\
\end{frame}

\begin{frame}{例題2}
  \begin{exampleblock}{問題}
    $d^{2}x/dt^{2} + 4x = e^{-t}$，（$t=0$ のとき $x=0,\ dx/dt=0$）を解け\\[2pt]
  \end{exampleblock}
\end{frame}

\begin{frame}[shrink=5]{例題2　解答}
  $\mathcal{L}[x(t)] = X(s)$ とおく．両辺をラプラス変換：\\[4pt]
  $(s^{2}+4)X(s) = \dfrac{1}{s+1}$ \quad $\therefore X(s) = \dfrac{1}{(s+1)(s^{2}+4)}$\\[4pt]
  部分分数分解：$\dfrac{1}{(s+1)(s^{2}+4)} = \dfrac{A}{s+1} + \dfrac{Bs+C}{s^{2}+4}$\\[4pt]
  通分：$1 = A(s^{2}+4) + (Bs+C)(s+1)$\\[4pt]
  $s=-1$：$A = \dfrac{1}{5}$，\quad $s^{2}$の係数：$B = -\dfrac{1}{5}$，\quad 定数項：$C = \dfrac{1}{5}$\\[4pt]
  $X(s) = \dfrac{1/5}{s+1} - \dfrac{1}{5}\cdot\dfrac{s}{s^{2}+4} + \dfrac{1}{10}\cdot\dfrac{2}{s^{2}+4}$
  \\[8pt]\textcolor{red}{\textbf{答}}：$x(t) = \dfrac{1}{5}\!\left[e^{-t} - \cos 2t + \dfrac{1}{2}\sin 2t\right]$
\end{frame}

\begin{frame}{演習問題}
  \begin{exampleblock}{自分で解いてみよう}
  \begin{description}
    \item[問2] 次の微分方程式を解け
    $(1)\ d^{2}x/dt^{2} - 5\,dx/dt + 6x = e^{t}$，（$t=0$ のとき $x=0,\ dx/dt=0$）\\
    $(2)\ d^{2}x/dt^{2} + 4\,dx/dt + 5x = 0$，（$t=0$ のとき $x=0,\ dx/dt=1$）\\
  \end{description}
  \end{exampleblock}
\end{frame}

\begin{frame}{問2　解答 (1)}
  (1) ラプラス変換：$(s^{2}-5s+6)X(s) = \dfrac{1}{s-1}$ \quad $\therefore X(s) = \dfrac{1}{(s-1)(s-2)(s-3)}$\\[4pt]
  部分分数分解：$s=1$：$A=\dfrac{1}{2}$，\quad $s=2$：$B=-1$，\quad $s=3$：$C=\dfrac{1}{2}$\\[4pt]
  $X(s) = \dfrac{1/2}{s-1} - \dfrac{1}{s-2} + \dfrac{1/2}{s-3}$\\[4pt]
  逆変換：$x(t) = \dfrac{1}{2}e^{t} - e^{2t} + \dfrac{1}{2}e^{3t}$
  \\[8pt]\textcolor{red}{\textbf{答}}：$x(t) = \dfrac{1}{2}e^{t} - e^{2t} + \dfrac{1}{2}e^{3t}$
\end{frame}

\begin{frame}{問2　解答 (2)}
  (2) ラプラス変換：$(s^{2}+4s+5)X(s) = 1$\\[4pt]
  $X(s) = \dfrac{1}{s^{2}+4s+5} = \dfrac{1}{(s+2)^{2}+1}$\\[4pt]
  逆変換：$x(t) = e^{-2t}\sin t$
  \\[8pt]\textcolor{red}{\textbf{答}}：$x(t) = e^{-2t}\sin t$
\end{frame}

\begin{frame}{例題3}
  \begin{exampleblock}{問題}
    $d^{2}x/dt^{2} - x = 0, x(0) = 0, x(1) = 1$ を解け（境界値問題）\\[2pt]
  \end{exampleblock}
\end{frame}

\begin{frame}{例題3　解答}
  $\mathcal{L}[x(t)] = X(s)$ とおく．$x(0)=0,\ x'(0)=\alpha$（未知）\\[4pt]
  原関数の微分法則：$\mathcal{L}[x''] = s^{2}X(s) - sx(0) - x'(0) = s^{2}X(s) - \alpha$\\[4pt]
  両辺をラプラス変換：$(s^{2}X(s)-\alpha) - X(s) = 0$\\[4pt]
  $(s^{2}-1)X(s) = \alpha$ \quad $\therefore X(s) = \dfrac{\alpha}{(s-1)(s+1)} = \dfrac{\alpha/2}{s-1} - \dfrac{\alpha/2}{s+1}$\\[4pt]
  逆変換：$x(t) = \dfrac{\alpha}{2}(e^{t} - e^{-t})$\\[4pt]
  境界条件 $x(1) = 1$ より $\dfrac{\alpha}{2}(e - e^{-1}) = 1$ \quad $\therefore \alpha = \dfrac{2}{e - e^{-1}}$
  \\[8pt]\textcolor{red}{\textbf{答}}：$x(t) = \dfrac{e^{t} - e^{-t}}{e - e^{-1}}$\\[4pt]
  {\footnotesize \textcolor{gray}{※ 境界値問題では未知の初期値 $\alpha$ を置き，もう一方の境界条件で決定する}}
\end{frame}

\begin{frame}{演習問題}
  \begin{exampleblock}{自分で解いてみよう}
  \begin{description}
    \item[問3] 次の微分方程式を解け（境界値問題）
    $(1)\ d^{2}x/dt^{2} + dx/dt = 0,\quad x(0)=2,\ x(1)=1+e^{-1}$\\
    $(2)\ d^{2}x/dt^{2} + x = 1,\quad x(0)=0,\ x(\pi/2)=0$\\
  \end{description}
  \end{exampleblock}
\end{frame}

\begin{frame}{問3　解答 (1)}
  (1) $x'(0)=\alpha$（未知）とおく．ラプラス変換：\\[4pt]
  $s(s+1)X(s) = 2s+\alpha+2$ \quad $\therefore X(s) = \dfrac{\alpha+2}{s} - \dfrac{\alpha}{s+1}$\\[4pt]
  逆変換：$x(t) = (\alpha+2) - \alpha e^{-t}$\\[4pt]
  境界条件 $x(1)=1+e^{-1}$ より $\alpha(1-e^{-1}) = -(1-e^{-1})$ \quad $\therefore \alpha=-1$
  \\[8pt]\textcolor{red}{\textbf{答}}：$x(t) = 1 + e^{-t}$
\end{frame}

\begin{frame}{問3　解答 (2) --- 1/2}
  (2) $x'(0)=\alpha$（未知）とおく．ラプラス変換：\\[4pt]
  $(s^{2}+1)X(s) = \dfrac{1}{s}+\alpha$ \quad $\therefore X(s) = \dfrac{1}{s(s^{2}+1)}+\dfrac{\alpha}{s^{2}+1}$\\[4pt]
  部分分数分解：$\dfrac{1}{s(s^{2}+1)} = \dfrac{A}{s} + \dfrac{Bs+C}{s^{2}+1}$\\[4pt]
  通分：$1 = A(s^{2}+1) + (Bs+C)s = (A+B)s^{2} + Cs + A$\\[4pt]
  定数項：$A=1$，\quad $s$の係数：$C=0$，\quad $s^{2}$の係数：$A+B=0 \Rightarrow B=-1$\\[4pt]
  $\therefore \dfrac{1}{s(s^{2}+1)} = \dfrac{1}{s} - \dfrac{s}{s^{2}+1}$
\end{frame}

\begin{frame}{問3　解答 (2) --- 2/2}
  $X(s) = \dfrac{1}{s}-\dfrac{s}{s^{2}+1}+\dfrac{\alpha}{s^{2}+1}$\\[4pt]
  逆変換：$x(t) = 1-\cos t+\alpha\sin t$\\[4pt]
  境界条件 $x(\pi/2)=0$ より $1+\alpha=0$ \quad $\therefore \alpha=-1$
  \\[8pt]\textcolor{red}{\textbf{答}}：$x(t) = 1-\cos t-\sin t$
\end{frame}

\begin{frame}{例題4}
  \begin{exampleblock}{問題}
    次の微分方程式の一般解を求めよ\\[2pt]
    $d^{2}x/dt^{2} - 4\,dx/dt + 4x = e^{2t}$
  \end{exampleblock}
\end{frame}

\begin{frame}{例題4　解答 (1/2)}
  $x(0)=\alpha,\ x'(0)=\beta$（未知）とおく．ラプラス変換：\\[4pt]
  $(s^{2}-4s+4)X(s) = \alpha s+\beta-4\alpha+\dfrac{1}{s-2}$\\[4pt]
  $(s-2)^{2}X(s) = \alpha(s-2)+(\beta-2\alpha)+\dfrac{1}{s-2}$\\[6pt]
  $\therefore\quad X(s) = \dfrac{\alpha}{s-2}+\dfrac{\beta-2\alpha}{(s-2)^{2}}+\dfrac{1}{(s-2)^{3}}$
\end{frame}

\begin{frame}{例題4　解答 (2/2)}
  $\mathcal{L}^{-1}\!\left[\dfrac{1}{(s-2)^{3}}\right]=\dfrac{t^{2}}{2}e^{2t}$ を用いて逆変換：\\[4pt]
  $x(t)=\alpha e^{2t}+(\beta-2\alpha)te^{2t}+\dfrac{t^{2}}{2}e^{2t}$\\[4pt]
  $C_{1}=\alpha,\ C_{2}=\beta-2\alpha$ を任意定数とする
  \\[8pt]\textcolor{red}{\textbf{答}}：$x(t)=C_{1}e^{2t}+C_{2}te^{2t}+\dfrac{t^{2}}{2}e^{2t}$（$C_{1},C_{2}$ は任意定数）
\end{frame}

\begin{frame}{演習問題}
  \begin{exampleblock}{自分で解いてみよう}
  \begin{description}
    \item[問4] 次の微分方程式の一般解を求めよ
    $(1) dx/dt - 2x = e^{3t}$\\
    $(2) d^{2}x/dt^{2} + 4x = \cos 2t$\\
  \end{description}
  \end{exampleblock}
\end{frame}

\begin{frame}{問4　解答 (1)}
  (1) $x(0) = \alpha$（未知）とおく．ラプラス変換：\\[4pt]
  $(s-2)X(s) = \alpha + \dfrac{1}{s-3}$ \quad $\therefore X(s) = \dfrac{\alpha}{s-2} + \dfrac{1}{(s-2)(s-3)}$\\[4pt]
  部分分数分解：$\dfrac{1}{(s-2)(s-3)} = \dfrac{-1}{s-2} + \dfrac{1}{s-3}$\\[4pt]
  $X(s) = \dfrac{\alpha-1}{s-2} + \dfrac{1}{s-3}$，\quad $C = \alpha-1$ を任意定数とする\\[4pt]
  逆変換：$x(t) = Ce^{2t} + e^{3t}$
  \\[8pt]\textcolor{red}{\textbf{答}}：$x(t) = Ce^{2t} + e^{3t}$（$C$ は任意定数）
\end{frame}

\begin{frame}{問4　解答 (2)}
  (2) $x(0)=\alpha,\ x'(0)=\beta$（未知）とおく．ラプラス変換：\\[4pt]
  $(s^{2}+4)X(s) = \alpha s + \beta + \dfrac{s}{s^{2}+4}$\\[4pt]
  $X(s) = \dfrac{\alpha s}{s^{2}+4} + \dfrac{\beta}{s^{2}+4} + \dfrac{s}{(s^{2}+4)^{2}}$\\[4pt]
  $\mathcal{L}^{-1}\!\left[\dfrac{s}{(s^{2}+4)^{2}}\right] = \dfrac{t}{4}\sin 2t$（像関数の微分から導出）\\[4pt]
  $C_{1} = \alpha,\ C_{2} = \beta/2$ を任意定数とする
  \\[8pt]\textcolor{red}{\textbf{答}}：$x(t) = C_{1}\cos 2t + C_{2}\sin 2t + \dfrac{t}{4}\sin 2t$
\end{frame}

\end{document}
